\begin{figure*}[ht!]
  \centering
    \figpanel{0.29\linewidth}{A}{fig3_panel_a_joint_distribution.pdf}\hfill
  \figpanel{0.34\linewidth}{B}{fig3_panel_b_nomenclature_examples.pdf}\hfill
  \figpanel{0.29\linewidth}{C}{fig3_panel_c_local_global_recovery.pdf}
  \vspace{0.5em}

  \figpanel{0.29\linewidth}{D}{fig3_panel_d_cap_joint_distribution.pdf}\hfill
  \figpanel{0.34\linewidth}{E}{fig3_panel_e_cap_nomenclature_examples.pdf}\hfill
  \figpanel{0.29\linewidth}{F}{fig3_panel_f_cap_local_global_recovery.pdf}
  \vspace{0.5em}

  \figpanel{0.30\linewidth}{G}{fig3_panel_g_cap_liver_reported_marker_recovery.pdf}\hfill
  \figpanel{0.30\linewidth}{H}{fig3_panel_h_cap_liver_de_stability.pdf}\hfill
  \figpanel{0.30\linewidth}{I}{fig3_panel_i_cap_liver_marker_retrieval.pdf}

  \caption{\textbf{Corpus-scale evaluation of marker genes and cell-type labels.} \textbf{A--C}. LLMarkers extracted corpus. \textbf{D--F}. CAP human marker profiles, restricted to cross-project comparisons. \textbf{G--I}. CAP liver expression-data example comparing a local myeloid dataset to the broader all-cells liver atlas. \textbf{A,D}. Label-marker agreement. Counts are cell-type pairs, grouped by reported label similarity and marker-gene set similarity. \textbf{B,E}. Inconsistent pairwise cell type labels. Swarm plots compare marker-gene set Jaccard for random different-label profile pairs and same-label profile pairs. Tables list different reported labels with high marker overlap. \textbf{C,F}. Globally identifiable markers. Each point is a recurrent label or ontology term. The $x$-axis is the mean fraction of markers that are unique to a profile within its source paper or CAP project. The $y$-axis is the mean fraction of those locally specific markers recovered anywhere in same-label or same-ontology-term profiles from other papers or projects. Immune labels are highlighted in orange. \textbf{G--I}. Reported markers were compared with one-vs-rest DE rankings in the local myeloid dataset and after mapping those labels to the broader all-cells liver context. Whiskers in \textbf{G} show 95\% intervals from 50 stratified cell-subsampling replicates.}
  \label{fig:cross_study_unification}%
  \claim{fig-cross-study-unification}{}%
  \source{analysis/build_fig3_local_global_marker_identifiability.py}%
  \source{analysis/build_cap_liver_local_global_deg.py}%
\end{figure*}
